% !TeX program = xelatex
\documentclass{UAmathtalk}
\usepackage{setspace}
\renewcommand{\when}{12:00 p.m.}
%\renewcommand{\where}{Hudson 0110. Zoom}
\renewcommand{\where}{Hudson 0110.\\ \href{https://albany.zoom.us/j/81219478289?pwd=N5fy7vxKBXqlQ4mAq4KaeaAM5iAHi4.1}{Zoom ID: 812\,1947\,8289 - Passcode: 233444}}
\author{Sushovan Majhi }
\address{George Washington University}
\urladdr{https://www.smajhi.com/}
\title{The Geometric Latschev’s Theorem: Euclidean Shape Reconstruction via~Vietoris–Rips~Shadow}
\date{Friday, April 17, 2026}


\begin{document}

\maketitle
%\renewcommand*{\abstractsize}{\normalsize}
\renewcommand*{\abstractsize}{\fontsize{13.5}{16.2}\selectfont}
\begin{abstract}
{The shadow of an abstract simplicial complex $\mathcal{K}$, whose vertices are in $\mathbb{R}^N$, is defined as the union of the convex hulls of its simplices. For a metric space $(S, d)$ at scale $\beta$, the Vietoris–Rips complex $\mathcal{R}_{\beta}(S)$ is the abstract simplicial complex where each $k$-simplex corresponds to $(k + 1)$ points in $S$ with a diameter at most $\beta$. Latschev’s theorem provides a qualitative guarantee for manifold reconstruction: for any closed Riemannian manifold $X$, there exists a scale $\epsilon_0 > 0$ such that for any $0 < \beta < \epsilon_0$, there is a $\delta > 0$ where any metric space $S$ within Gromov–Hausdorff distance $\delta$ of $X$ yields a Vietoris–Rips complex homotopy equivalent to $X$. Recently, Latschev’s theorem has been quantified, allowing $X$ to be a more general geodesic space. When $X \subset \mathbb{R}^N$ is a Euclidean geodesic space (e.g., submanifold, graph), we address the theorem’s geometric analog: under what conditions is the shadow of the Vietoris–Rips complex of a Hausdorff-close Euclidean sample $S \subset \mathbb{R}^N$ both homotopy equivalent and Hausdorff-close to $X$? Unlike the abstract complex, the shadow provides a geometric embedding within the host space, which is essential for the practical reconstruction of Euclidean shapes. In this talk, we discuss recent developments in answering this question and explore their implications for faithful reconstruction of low-dimensional submanifolds and Euclidean-embedded graphs.}
\end{abstract}
\end{document}

